\section{引言}

本文的大多数具体结果是用格点场论方法得到的。这些结果构成一个更大、相当雄心勃勃且颇为思辨性的纲领的第一部分；该纲领是在连续统中定义的，不依赖于任何正规化方案。本节的主要内容就是勾画这一纲领，其目的在于为后面各节给出的格点结果提供恰当的视角。在整篇文章中，我们会不时回到连续统的语言，以强调在我们看来哪些结构是连续统的性质。


纯四维 $SU(N)$ 规范理论在大距离处强相互作用，在短距离尺度上弱相互作用。在 $SU(3)$ 中，这两种区域之间的渡越区（cross-over）相对较窄，没有好的解析近似方法适用于它；我们也没有在强耦合下做计算的连续统方法。随着色数 $N$ 的增大，渡越区变窄。

此处，作为对某些早期工作 \cite{old1, gw12} 思路的直接延续，我们报告一条改进了的进攻路线，并沿此方向迈出几步。旧的想法是：瞬子表明，在无穷大 $N$ 极限下，在一个被明确界定的特定标度处存在一个相变，弱耦合与强耦合物理可以在那一点衔接起来。这个想法从未被成功地实现过，因为其中太多部分要么错误要么缺失。在此期间人们学到了许多东西，我们愿意再次尝试。现在让我们简要概述我们做法中的新要素：

虽然``渡越区在无穷大 $N$ 时坍缩为一点并成为一个相变''这一猜想由来已久，但我们现在明白，相变发生的标度并不是一个同时适用于所有可观测量、处处成立的普适量；它也不必然影响任何局域可观测量：例如它在自由能密度中就看不到。最有趣的大 $N$ 相变未必是体（bulk）相变。

近来关于大 $N$ 的格点研究在欧氏环面上 QCD 的平面极限中于整体临界标度处产生了这样的体相变，但相关的相变把系统带出了零温、无限体积的区域，并且依赖时空的整体拓扑。这无法直接解释例如少数粒子的高能过程如何强子化。

近来的格点工作还表明，特定的非局域可观测量在被胀大时会经历大 $N$ 相变 \cite{bulk2,knn}，而这发生的同时欧氏时空体积在所有方向上都保持无限。临界标度依赖于可观测量的选择；可以想象给这些可观测量附上一整族胀大的像，而大 $N$ 相变发生在膨胀轴上的某处。

如果存在平面 QCD 的弦等价描述，人们会猜测膨胀轴以某种方式与第五个坐标相关联，正如经常被假定的那样 \cite{poly3,malda4}。一个有意义的高维表述预设了理论在延拓空间中的局域定义；在本文末尾我们将对这一思辨补充一点看法，但它会是含糊的。

我们为旧纲领添加的另一新要素是对矩阵模型中某些大 $N$ 相变之普适特征的现代理解。尽管表面上看，在与对称群相联系的指标空间中没有局域性的概念，但这些模型对二指标表示容许按二维曲面亏格分类的 't Hooft 双线展开。一种二维局域性概念适用于这些曲面上的物理，而相关的 Liouville 型模型代表与大 $N$ 相变相伴的、稳定性程度不同但有限的若干普适类 \cite{doublesca5}。这些模型大多精确可解。

我们希望将要展示的相变属于某个这样的普适类。我们对相关普适类有一个猜想，而且它有一个精确可解的矩阵模型代表元。这意味着，在有限但足够大的 $N$ 下，当相变弥散成渡越区时，对精心挑选的、经历渡越的可观测量存在一种描述：用普适标度函数来描写，而这些普适标度函数只依赖少量非普适参数；这些参数通常作为确定已知普适标度函数之方程的边界条件出现。

正是通过这些参数，渡越区才与两侧的弱耦合和强耦合区域连接起来。四维膨胀参数在其临界值附近扮演着一个耦合常数的角色；在大 $N$ 相变的相应矩阵模型描述中，该耦合变为临界的。借助（双重）标度极限机制，膨胀参数与二维物理中的一个标度联系起来，后者编纂了大 $N$ 相变的普适特征。这一联系将涉及一个非平凡的幂次以及一个相关的有量纲参数。

弱耦合与强耦合区域各自受制于自己的普适行为：弱耦合侧是场论的，强耦合侧则是（我们猜测）有效弦式的。于是，每种普适描述——场论的、大 $N$ 的和弦式的——所需的自由非普适参数，在算符被胀大时彼此匹配；几种普适性机制共存于同一系统中，分别控制不同的标度区间。

若上述情景在 $N=\infty$ 成立，我们可以把它应用于纯 $SU(3)$ 规范理论，并有望给出用 $\Lambda_{\rm QCD}$ 表达弦张力的计算，精度达到 $\frac{1}{ N^2}\approx 10\%$ 的量级。对一个物理学家来说，这可能比禁闭的数学证明更有价值。我们距离实现这一目标还很远。

要克服的第一道障碍是找到好的可观测量；我们主张，正确的问题不是一般性地问是否存在大 $N$ 相变，而是能否找到既经历这种相变、又能在极短和极长距离尺度上分别（由不同理论）加以控制的好可观测量。局域可观测量不太可能经历这种相变。因此，四维欧氏动量空间中备受珍视的解析性质得以保全，物理的场论描述的基本幺正性与协变性得以维持。

本文我们在格点上定义了一类涂抹（smeared）的威尔逊圈可观测量，意图使它们具有有限的连续统极限并满足如下性质：
\begin{itemize}
\item 把绕圈的平行输运看作由一个 $SU(N)$ 矩阵来实现是正确的。这一陈述必须在理论和算符的重整化完成之后仍然成立。
\item 当我们对圈作标度变换时，若 $N=\infty$，其本征值分布中会发生一次相变。
\item 涂抹圈算符在被缩小时可用微扰论计算。
\item 其形状依赖在圈被放大时由一个有效弦理论很好地描述。
\item 有可能通过连续统的非微扰方法（例如瞬子）看出，当圈从非常小被胀大时是如何逼近相变的，从而可以指望计算渡越区弱耦合端的匹配。
\end{itemize}
在大尺度一端还需要有办法计算一次匹配；据推测这之所以可行，是因为渡越区的大尺度端具有某种弦描述（来自相应的普适矩阵模型），它可以与正确的四维有效弦理论光滑地拼接起来。那个有效弦理论例如可以类似于 \cite{ps6}。这一点我们尚未确立，留待将来工作。尽管如此，基于已有的结果，我们相信自己已经找到了足够``好''的非局域可观测量。

本质上说，我们所勾勒的计算是把有效场论 \cite{effft7} 计算的逻辑扩展到纯规范 QCD 强—弱渡越区的一种尝试。此外，该纲领与旧的 CDG 瞬子方法 \cite{cdg75} 精神相同：他们曾试图用瞬子气体方法来确定强子的有效袋模型描述中的参数。

如果我们的情景对纯胶子系统成立，那么就有可能在稍后阶段重新引入费米子，至少作为没有反作用的探针，并看看能否构造出同样经历大 $N$ 相变的、含费米子场的强子可观测量。特别有趣的是为自发手征对称性破缺找到一个类似的计算纲领。可以用严格手征费米子~\cite{overlap} 对此进行数值检验。

下一节中我们在格点上引入我们的可观测量，并给出数值结果，以支持如下主张：我们的可观测量具有连续统极限，并且在胀大下经历大 $N$ 相变。在随后几节中，我们讨论格点构造的连续统观点以及我们对大 $N$ 普适类的猜想。叙述不可避免地将越来越具思辨性。最终，我们希望把格点模拟的作用缩减为严格的定性支撑，即为控制着我们计算需要衔接的各种区域的诸普适类假设提供支持。计算本身将完全是解析的。不言而喻，任何解析步骤都需要对照数值工作加以检验，因为这是一片未经勘测的领域。
